% Generated draft: requires mathematical and source review.
\documentclass[11pt]{article}
\usepackage[margin=25mm]{geometry}
\usepackage{fontspec}
\setmainfont{Latin Modern Roman}
\usepackage{amsmath,amssymb,mathtools,unicode-math}
\setmathfont{Latin Modern Math}
\usepackage{enumitem,needspace,xurl,hyperref}
\setlength{\parindent}{0pt}
\setlength{\parskip}{5pt}
\setlength{\emergencystretch}{3em}
\newcommand{\N}{\mathbb{N}}
\newcommand{\Z}{\mathbb{Z}}
\newcommand{\Q}{\mathbb{Q}}
\newcommand{\R}{\mathbb{R}}
\newcommand{\C}{\mathbb{C}}
\newcommand{\F}{\mathbb{F}}
\newcommand{\A}{\mathbb{A}}
\newcommand{\PP}{\mathbb{P}}
\newcommand{\E}{\mathbb{E}}
\newcommand{\eps}{\varepsilon}
\newcommand{\dd}{\,\mathrm{d}}
\begin{document}
\section*{UnsolvedMath ID: 20002069}
\textbf{Source problem code:} catalogue entry
\subsection{Definitions and mathematical statement}
Question 1.11 asks how symplectic field theory (SFT) can be extended so that its invariants can be computed by cutting along successively lower-dimensional strata. Its TQFT analogy assigns numbers to closed \(n\)-manifolds, vector spaces to \((n-1)\)-manifolds, categories to \((n-2)\)-manifolds, then higher categories in lower dimensions, with compatible gluing axioms. It supplies no precise geometric bordism category with corners, tangential/contact decorations, target higher category, or depth of extension. The request is a construction program rather than a single fully specified existence theorem.

The aid gives a formal categorical lift under explicit algebraic hypotheses. Let \(\Bbbk\) be a field and let a strict symmetric monoidal cobordism functor assign unital differential graded \(\Bbbk\)-algebras and algebra maps, with variance fixed by that functor. For an algebra \(A\), let \(\operatorname{Perf}(A)\) be its idempotent-complete category of perfect right modules. An algebra map \(f:A\to B\) induces
\[
 f_!:\operatorname{Perf}(A)\longrightarrow\operatorname{Perf}(B),
 \qquad M\longmapsto M\otimes_A^{\mathbf L}B.
\]
Derived tensor associativity gives the coherent composition identification for \(A\to B\to C\), and tensor product over \(\Bbbk\) gives the monoidal structure in the standard derived Morita setting. The aid calls this the canonical extension-of-scalars Morita lift. Its application to Pardon's rigorous contact-homology cobordism functor is a categorical shadow; the exact functor model, coefficients, variance, and coherence hypotheses for that application are not reproduced.

This lift reprocesses already assigned algebras and cobordism maps. It does not by itself assign invariants to new codimension-two contact-geometric strata or prove geometric cutting/gluing with corners. It also concerns contact homology, not an asserted complete analytic construction of every SFT operation.

The aid separately gives a properness obstruction to one route to a fully extended theory. For \(T^*S^1\), let \(L\) be a cotangent fiber. It reports
\[
 HW^*(L,L;\Bbbk)\cong\Bbbk[\Z]
 =\Bbbk[t,t^{-1}],
\]
with its infinite-dimensional degree-zero cohomology. A small \(\Bbbk\)-linear category is proper when every morphism complex has finite-dimensional total cohomology. Hence the wrapped Fukaya category \(\mathcal W(T^*S^1)\) is not proper. In the standard Morita higher-category setting where full dualizability requires smoothness and properness, it therefore cannot automatically be a fully dualizable point value of a fully extended TQFT.

The target qualification matters: this is not a theorem excluding all extended SFT constructions, other categorical targets, or partial extensions only to codimension two. The aid concludes only that sectorial descent and stop removal, together with wrapped categories, do not alone provide the requested fully extended standard-Morita theory. They do not remove the geometric, analytic, coherence, and dualizability tasks.
\subsection{Short English statement}
The source asks for SFT with invariants on lower-dimensional strata and compatible cutting rules. The aid lifts an existing algebra-valued contact-homology functor to perfect-module categories, but that adds no new geometric strata. Its cotangent-circle example obstructs using arbitrary wrapped categories as automatically fully dualizable point values in the standard Morita target.
\subsection{Sources}
\begin{itemize}
\item Local UnsolvedMath registry, permanent ID 20002069, exact source code catalogue entry; canonical \texttt{aim-geometry-notes.json}, \texttt{notes[406]}, tag \texttt{question}; AIM Workshop Problem Lists, Geometry, \emph{Holomorphic curves in contact geometry}, Question 1.11. Created and updated 14 August 2026; statement status \texttt{reconstructed\_unverified}.
\item AIM source version 15 October 2003, according to supplied verification; \url{https://aimath.org/WWN/contactgeom2/contactgeom2.pdf}. The verification calls its transcription a conservative reconstruction. The PDF was not accessed.
\item The aid names Pardon's contact-homology functor and invokes derived Morita theory, the cotangent-fiber wrapped calculation, sectorial descent, and stop removal, without full publication entries, theorem numbers, or separate URLs. No authoritative exact theorem citation is supplied for the proposed application or dualizability target.
\item Attempt 1: \texttt{partial\_result}, problem status at run \texttt{open}, rigor label \texttt{proved}, novelty confidence \texttt{low}; classification \texttt{PARTIAL-PROGRESS}, registry status \texttt{partially\_solved}. Candidate: formal Morita lift and a standard-target properness obstruction. These are dated supplied assessments, not acceptance.
\item Ulam AI / UnsolvedMath Contributors; CC BY 4.0. Machine-generated research aid, not peer reviewed; independently verify before citation.
\end{itemize}
\end{document}
